% !TeX root = ../../main.tex
\subsection{疑問詞とbe動詞を含む文の疑問文}
    疑問詞とbe動詞を含む文の疑問文について考える前に、何点か大切なことを復習しておこう。

    \begin{theorembox}{\textbf{be動詞の用法}}
        \begin{enumerate}
            \item 英語の文において動詞は必ず\fitblankbf[]{1つ}しか使えない
            \item be動詞を含む文の疑問文は、\fitblankbf[]{be動詞}を\fitblankbf[]{文頭}に出す
            \item be動詞を含む文の否定文は、\fitblankbf[]{be動詞}の直後に\fitblankbf[]{not}をつける
        \end{enumerate}
    \end{theorembox}

    まずはbe動詞を含む文の用法から復習しておこう。

    \begin{itemize}
        \item This is \underline{your bag}.
        \par\vspace{0.3em}
        \noindent\makebox[12em][l]{疑問文へ}\ $\longrightarrow$\ \dialogueblankinline{Is this your bag?}
        \par\vspace{0.5em}
        \noindent\makebox[12em][l]{否定文へ}\ $\longrightarrow$\ \dialogueblankinline{This is not your bag.}
        \par\vspace{0.5em}
        \noindent\makebox[12em][l]{下線部を問う疑問文へ}\ $\longrightarrow$\ \dialogueblankinline{What is this?}
        \par\vspace{0.5em}
        \noindent\makebox[12em][l]{下線部の\uwave{場所}を問う疑問文へ}\ $\longrightarrow$\ \dialogueblankinline{Where is this?}
    \end{itemize}

\subsection{疑問詞と一般動詞を含む文の疑問文}
    疑問詞と一般動詞を含む文の疑問文について考える前に、be動詞のとき同様に何点か大切なことを復習しておこう。

    \begin{theorembox}{\textbf{一般動詞の用法}}
        \begin{enumerate}
            \item 英語の文において動詞は必ず\fitblankbf[]{1つ}しか使えない
            \item 一般動詞を含む文の疑問文は、基本的には\fitblankbf[]{Do}を\fitblankbf[]{文頭}に出す
            \item 一般動詞を含む文の否定文は、基本的には\fitblankbf[]{don't}と使う
            \item 主語が三人称、単数、時制が現在のときは\fitblankbf[]{Does}、\fitblankbf[]{doesn't}を使う
            \item 時制が過去のときは、\fitblankbf[]{Did}、\fitblankbf[]{didn't}を使う
        \end{enumerate}
    \end{theorembox}

    まずは一般動詞を含む文の用法から復習しておこう。

    \begin{itemize}
        \item You play \underline{soccer} every Sunday.
        \par\vspace{0.3em}
        \noindent\makebox[12em][l]{疑問文へ}\ $\longrightarrow$\ \dialogueblankinline{Do you play soccer every Sunday?}
        \par\vspace{0.5em}
        \noindent\makebox[12em][l]{否定文へ}\ $\longrightarrow$\ \dialogueblankinline{You don't play soccer every Sunday.}
        \par\vspace{0.5em}
        \noindent\makebox[12em][l]{下線部を問う疑問文へ}\ $\longrightarrow$\ \dialogueblankinline{What do you play every Sunday?}
        \par\vspace{0.5em}
        \noindent\makebox[12em][l]{下線部の\uwave{場所}を問う疑問文へ}\ $\longrightarrow$\ \dialogueblankinline{Where do you play soccer every Sunday?}
    \end{itemize}
\newpage

\subsection{疑問詞と助動詞を含む文の疑問文}
    前に倣って、疑問詞と助動詞を含む文の疑問文について考える前に、何点か大切なことを復習しておこう。

    \begin{theorembox}{\textbf{助動詞の用法}}
        \begin{enumerate}
            \item 助動詞の後ろは必ず\fitblankbf[]{動詞の原形}
            \item 助動詞を含む文の疑問文は、\fitblankbf[]{助動詞}を\fitblankbf[]{文頭}に出す
            \item 助動詞を含む文の否定文は、\fitblankbf[]{助動詞}の直後に\fitblankbf[]{not}をつける

            したがって、助動詞を含む文の用法は \fitblankbf[]{be動詞}を含む文の用法と同じである。
        \end{enumerate}
    \end{theorembox}

    まずは助動詞を含む文の用法から復習しておこう。

    \begin{itemize}
        \item He can play \underline{the guitar} very well.
        \par\vspace{0.3em}
        \noindent\makebox[12em][l]{疑問文へ}\ $\longrightarrow$\ \dialogueblankinline{Can he play the guitar very well?}
        \par\vspace{0.5em}
        \noindent\makebox[12em][l]{否定文へ}\ $\longrightarrow$\ \dialogueblankinline{He can't play the guitar very well.}
        \par\vspace{0.5em}
        \noindent\makebox[12em][l]{下線部を問う疑問文へ}\ $\longrightarrow$\ \dialogueblankinline{What can he play very well?}
    \end{itemize}
\vspace{1.5em}

    これで、疑問詞を使う疑問文を使いこなすための最低限の知識は揃ったことになる。ここからは応用力で乗り越えてみてほしい。

    \begin{itemize}
        \item Her name is \underline{Hana}.
        \par\vspace{0.3em}
        \noindent\makebox[12em][l]{下線部を問う疑問文へ}\ $\longrightarrow$\ \dialogueblankinline{What is her name?}
        \par\vspace{0.5em}
        \item Tom walks to school \underline{with his brother}.
        \par\vspace{0.3em}
        \noindent\makebox[12em][l]{下線部を問う疑問文へ}\ $\longrightarrow$\ \dialogueblankinline{Who does Tom walk to school with?}
        \par\vspace{0.5em}
        \item They studied English \underline{at school}.
        \par\vspace{0.3em}
        \noindent\makebox[12em][l]{下線部を問う疑問文へ}\ $\longrightarrow$\ \dialogueblankinline{Where did they study English?}
        \par\vspace{0.5em}
        \item You must finish your homework \underline{by 8:00}.
        \par\vspace{0.3em}
        \noindent\makebox[12em][l]{下線部を問う疑問文へ}\ $\longrightarrow$\ \dialogueblankinline{When must you finish your homework?}
        \par\vspace{0.5em}
        \item I should buy \underline{this dress}.
        \par\vspace{0.3em}
        \noindent\makebox[12em][l]{下線部を問う疑問文へ}\ $\longrightarrow$\ \dialogueblankinline{What should I buy?}
        \par\vspace{0.5em}
    \end{itemize}

\vspace{2em}

    疑問詞を使うときだけでなく、英語を学習する上で常に気にしておかなければならないことがある。それは、大きく \fitblankbf[]{動詞}である。
    主語がない文は\uwave{命令文}として存在するが、\textbf{動詞のない文は存在しない}。したがって、動詞が英文の核ということである。

    具体的にいうと、前述の通り、\fitblankbf[]{使う動詞}によって疑問文の形が変わる。そのため、\uwave{使いたい動詞を先に決めてから}
    疑問詞を使えるようになると、大したもんである。

\newpage

